% ============================================================
% REMERCIEMENTS
% ============================================================
\cleardoublepage
\thispagestyle{empty}

\vspace*{3cm}

\begin{center}
\Large \textbf{REMERCIEMENTS}
\end{center}

\vspace{1cm}

Ce mémoire est l'aboutissement d'un parcours de recherche rendu possible grâce au soutien de plusieurs acteurs que je tiens à remercier chaleureusement.

Mes remerciements s'adressent tout d'abord à mon encadrant, \textbf{ Monsieur Gracieux HOUNNA}, pour avoir cru en la faisabilité du pipeline \AZIZA\ et pour m'avoir accordé la liberté scientifique d'explorer l'hybridation entre les algorithmes NLP classiques et les grands modèles de langage. Votre encadrement et vos conseils avisés ont été précieux tout au long de mes travaux.

Je témoigne ma profonde gratitude au Directeur de l'école, \textbf{ Monsieur Gaston EDAH}, ainsi qu'à mon professeur principal, \textbf{ le Dr Ratheil HOUNDJI}, pour la rigueur de leur suivi, la qualité de leur encadrement pédagogique et leur soutien constant durant mon parcours.

J'adresse mes vifs remerciements aux experts et enseignants de français qui ont accepté de prêter leur expertise pour l'annotation de notre corpus de test, permettant ainsi de calibrer scientifiquement notre application.

Que les membres du jury trouvent ici l'expression de mon profond respect pour l'honneur qu'ils me font en évaluant ce travail et pour le temps consacré à l'examen de ces travaux.

Je tiens également à saluer mes camarades de classe, avec qui j'ai partagé d'intenses moments de codage, de précieux partages techniques et des défis stimulants qui m'ont permis de m'améliorer au quotidien.

Enfin, je remercie sincèrement mes amis ainsi que le reste de ma famille pour leur présence indéfectible, leur soutien moral et leurs encouragements constants qui ont constitué un havre de motivation tout au long de cette aventure académique.

\vfill
